% !TeX root = ../main.tex
\chapter{Singular directions and two-point functions}\label{f:microlocal}

A delta function on a sheet is singular at every point of the sheet. Yet a
probe moving across the sheet can register a perfectly well-defined impulse,
whereas a probe confined to the sheet would be asked to sample that impulse
at zero distance for its entire history. Singular support alone does not
distinguish the two measurements. We need to know in which directions a
localized field carries arbitrarily high frequencies.

\section{A Fourier test with two localizations}
Let $u$ be a distribution on an open subset of $\mathbb R^n$. Localize it
first in position: take a smooth compactly supported cutoff $\chi$ equal to
one on a neighborhood of $x_0$. Then
\[
  \widehat{\chi u}(\xi)=\langle u,\chi(x)e^{-ix\cdot\xi}\rangle
\]
is a smooth function of at most polynomial growth. Indeed the test function
in this pairing has fixed compact support, and differentiating it a finite
number of times in $x$ costs at most a polynomial in $\xi$. If $u$ is
smooth on the support of $\chi$, integration by parts gives rapid decay in
all directions. Conversely, rapid decay in all directions allows Fourier
inversion after multiplication by every monomial $(i\xi)^\alpha$; hence
$\chi u$ is smooth. The second localization keeps only directions in an
open cone $\Gamma$: if $\xi\in\Gamma$, so is $a\xi$ for every $a>0$.
One ray by itself is not an open cone in dimension greater than one.

\begin{definition}\label{f:wf-test}
A pair $(x_0,\xi_0)$, $\xi_0\ne0$, is \emph{regular for $u$} if there
exist such a cutoff $\chi$ and an open cone $\Gamma$ containing $\xi_0$
for which, for every nonnegative integer $N$,
\begin{equation}
 \sup_{\xi\in\Gamma}(1+|\xi|)^N|\widehat{\chi u}(\xi)|<\infty .
 \label{f:wf-decay}
\end{equation}
The complementary pairs form the \emph{wavefront set} $\operatorname{WF}(u)$.
It is closed in the cotangent bundle with the zero section removed, and
is invariant under positive rescaling of the covector.
\end{definition}

The same cutoff and cone work for nearby pairs, which explains closedness.
A more useful fact is that one can shrink the position neighborhood even
when this changes the Fourier transform everywhere. If $\psi$ is supported
where $\chi=1$, the transform of $\psi u$ is the convolution of
$\widehat\psi$ with $\widehat{\chi u}$, divided by $(2\pi)^n$.
Take a smaller cone whose unit directions have closure inside $\Gamma$.
For $\xi$ in this cone, split the convolution integral over $\eta$ into
$\eta\in\Gamma$ and its complement. In the first part use
$1+|\xi|\le(1+|\xi-\eta|)(1+|\eta|)$: every power of the second factor
can be absorbed by \eqref{f:wf-decay}, and $\widehat\psi$ is rapidly
decreasing. In the second part angular separation gives
$|\xi-\eta|\ge c(|\xi|+|\eta|)$ for some $c>0$. There the decay of
$\widehat\psi(\xi-\eta)$ absorbs both the polynomial growth of
$\widehat{\chi u}(\eta)$ and any desired power of $|\xi|$.

\begin{proposition}
The projection of $\operatorname{WF}(u)$ onto the base is the ordinary
singular support of $u$.
\end{proposition}
\begin{proof}
A smooth neighborhood supplies a single cutoff good in all directions.
Conversely, if no direction over $x_0$ is singular, cover the unit sphere
by finitely many of the \emph{smaller} cones just described. Choose one
cutoff equal to one near $x_0$ and supported where all the corresponding
position cutoffs are one. Its transform is rapidly decreasing on each of
the finitely many cones, hence everywhere. Fourier inversion now gives a
smooth representative near $x_0$.
\end{proof}

\section{A sheet and a one-sided singularity}
In coordinates $(y,s)\in\mathbb R^{n-1}\times\mathbb R$, the distribution
$\delta(s)$ acts by integrating a test function over $s=0$. With dual
coordinates $(\eta,\tau)$ one has
\[
 \widehat{\chi\delta(s)}(\eta,\tau)
    =\int e^{-iy\cdot\eta}\chi(y,0)\,\mathrm{d}y .
\]
In a cone with $|\eta|\ge c|(\eta,\tau)|$, rapid decay follows by
integrating in $y$. In contrast, when $\chi$ equals one near a point of
the sheet its restriction to the sheet is not identically zero. Its
transform is nonzero at some fixed $\eta_*$; the sequence
$(\eta_*,\tau)$ tends in direction to either normal ray as
$\tau\to\pm\infty$, without any decay. Thus the wavefront set consists
precisely of $(y,0;0,\tau)$ with $\tau\ne0$. The two signs matter.

A distribution can have just one sign even at a point singularity.
Set $b_\pm(x)=(x\pm i0)^{-1}$, the limits in distributions of
$(x\pm i\epsilon)^{-1}$ as $\epsilon\downarrow0$. Their real part tends
to the principal value of $1/x$; their imaginary parts tend to
$\mp\pi\delta_0$, since $\epsilon/(x^2+\epsilon^2)$ has integral $\pi$
and concentrates at zero. To find the singular directions, use the actual
Fourier sign, not the sign of the imaginary part. For instance
\[
  (x+i\epsilon)^{-1}=-i\int_0^\infty e^{ipx-\epsilon p}\,\mathrm{d}p.
\]
The conjugate formula has $+i$ and $e^{-ipx}$. With forward phase
$e^{-ix\xi}$ and inverse factor $1/(2\pi)$ this gives
\begin{equation}
 \widehat b_+(\xi)=-2\pi i\,\mathbf 1_{\xi>0},\qquad
 \widehat b_-(\xi)= 2\pi i\,\mathbf 1_{\xi<0} .
 \label{f:wf-bv}
\end{equation}
These are transforms of the \emph{unlocalized} distributions. To perform
the defining test, multiply by $\chi$ and convolve. The answer for $b_+$
is $-i\int_{-\infty}^\xi\widehat\chi(v)\,\mathrm{d}v$. It decreases
rapidly toward $-\infty$ but tends to $-2\pi i\chi(0)$ toward $+\infty$.
The answer for $b_-$ is $i\int_\xi^\infty\widehat\chi(v)\,\mathrm{d}v$.
Consequently $\operatorname{WF}(b_\pm)=\{(0,\xi):\ \pm\xi>0\}$. Although
both singular supports are $\{0\}$, the singular frequencies are disjoint.

\begin{exercise}
Determine the two limiting high-frequency values of
$\widehat{\chi\,\mathrm{pv}(1/x)}$ when $\chi=1$ near zero. Do the same
for $\chi\delta_0$. Why is the singularity of neither distribution
one-sided?
\end{exercise}
\noindent\emph{Hint.} Take the half-sum of $b_+$ and $b_-$ in the first
case and $(b_--b_+)/(2\pi i)$ in the second. On the positive ray the
answers tend to $-i\pi$ and $1$, respectively; now reverse the ray.

\section{Restriction and multiplication as pullbacks}
For a smooth map $F:Y\to X$, the transpose of the differential sends a
covector $k$ at $F(y)$ to $dF_y^T k$ at $y$. A rapid phase on $X$ may
cease to oscillate after composition with $F$: its new differential can
be zero. Define the set of offending directions of the map by
\[
 N_F=\{(F(y),k):k\ne0,\ dF_y^Tk=0\}.
\]
For an embedding this is its nonzero conormal bundle: its covectors
annihilate the tangent space of the embedded submanifold.

\begin{theorem}[H\"ormander's pullback theorem, smooth manifolds \cite{HormI}]
\label{f:wf-pullback}
If $u$ is a distribution on $X$ and $N_F\cap\operatorname{WF}(u)$ is
empty, the distributional pullback $F^*u$ exists canonically, agrees
with composition on smooth functions, and satisfies
\[
 \operatorname{WF}(F^*u)\subset
 \{(y,dF_y^Tk):(F(y),k)\in\operatorname{WF}(u)\}.
\]
Here distributions are scalar, acting on test densities. For a
diffeomorphism the inclusion is equality, by applying the theorem to
the inverse.
\end{theorem}

The transformation law also explains why these directions are covectors
rather than vectors. In coordinates the oscillatory factor before Fourier
transformation on $Y$ is $e^{i F(y)\cdot k}$; its new Fourier phase has
differential $dF_y^Tk-\eta$. Directions separated from the possible
stationary ones can be estimated by integration by parts. In particular,
straightening a smooth sheet by a diffeomorphism turns the calculation for
$\delta(s)$ into the intrinsic statement: a nonvanishing smooth density
of delta mass on a sheet has exactly its nonzero conormals as wavefront
set. Smooth nonzero density factors do not change this conclusion.

For multiplication use the diagonal map $d:x\mapsto(x,x)$ and set
$uv=d^*(u\otimes v)$. Its transpose sends $(k,l)$ to $k+l$.
The tensor-product wavefront estimate allows a singular covector from
either factor alone (zero in the other slot), or one from each. Thus a
sufficient test is that no $(x,k)\in\operatorname{WF}(u)$ have
$(x,-k)\in\operatorname{WF}(v)$. In particular $b_+^2$ passes, whereas
$b_+b_-$ does not. The latter's common-width regularizations give
$1/(x^2+\epsilon^2)$, with pairing asymptotic to
$\pi\psi(0)/\epsilon$ on a test $\psi$ with $\psi(0)\ne0$. A failed
sufficient test is not a general impossibility theorem: the square of a
bounded step function, for example, is defined by ordinary almost
everywhere multiplication.

For a kernel $K(x,y)$ it is useful to record $(x,k;y,l)$ in its
\emph{primed} relation when $(x,k;y,-l)\in\operatorname{WF}(K)$. Then the
delta kernel of the identity relates a covector to itself. There is a
simple composition rule when every singular pair of both kernels has two
nonzero entries. For $K(x,y),L(y,z)$ assume also that the projection of
their joint support in $(x,y,z)$ onto $(x,z)$ is proper (compact sets
have compact inverse images there). Multiply their lifts to the triple
space and integrate in $y$ with the chosen smooth density. Cancellation in
the product test cannot occur: a singularity from $K$ has nonzero $x$
covector, one from $L$ has nonzero $z$ covector. Under integration only
singular directions with zero \emph{total} $y$ covector can remain. The
primed relation of the result is therefore contained in the composition
of the two primed relations \cite{HormI}. Dropping the support assumption
can instead leave a divergent integral even for smooth kernels.

\begin{exercise}
Pull the sheet $\delta(s)$ in the $(y,s)$ plane back by the curves
$v\mapsto(v,av)$ for $a\ne0$ and $v\mapsto(v,v^2)$. Where does the test
fail? For the second curve, try a smooth approximation
$\delta_\epsilon(s)=\epsilon^{-1}\rho(s/\epsilon)$ with $\rho\ge0$,
$\rho(0)>0$, compact support and unit integral.
\end{exercise}
\noindent\emph{Solution.} The transverse pullback is
$|a|^{-1}\delta(v)$. At the tangency the conormal $\mathrm{d}s$ pulls
back to zero. Setting $v=\sqrt\epsilon\,w$ in the regularized pairing
gives leading term $\epsilon^{-1/2}\psi(0)\int\rho(w^2)\,\mathrm{d}w$.
The transverse formula has no limiting analogue here.

\section{Which singularities can a wave transport?}
Now take $P=\Box_g-m^2-\xi R$ for a real scalar on a smooth Lorentzian
spacetime, with $\Box_g=\nabla^a\nabla_a$, $R$ the scalar curvature and
$\xi$ a constant curvature coupling. Its characteristic covectors are the
nonzero zeros of
$p(x,k)=g^{ab}k_a k_b/2$. With Fourier symbol $\partial_a\mapsto ik_a$,
$p$ is the principal symbol multiplied by $-1/2$; this rescaling only
reparametrizes its Hamilton curves. A hypersurface is
\emph{non-characteristic for $P$} if its nonzero conormals are not
characteristic. This is a property of an operator and a surface, not a
test for restricting an arbitrary distribution to the surface. For
example, $t=0$ in Minkowski space is non-characteristic for $P$, yet
$\delta(t)$ has precisely the conormals that obstruct its own trace.

\begin{theorem}[Microlocal elliptic regularity and propagation, scalar wave operator \cite{HormIII}]
\label{f:wf-propagation}
Let $u$ be a distribution in a smooth boundaryless open region, and
let $O$ be an open conic set in its nonzero cotangent bundle on which
$Pu$ is microlocally smooth. Then $\operatorname{WF}(u)\cap O$ lies in
$p=0$ and is a union of maximal segments in $O$ of the Hamilton curves
of $p$ there.
\end{theorem}

The hypothesis on $Pu$ says $O\cap\operatorname{WF}(Pu)=\varnothing$.
The Hamilton curves in the characteristic set are called
\emph{bicharacteristics}. The real principal type hypothesis of the
propagation theorem holds here: $\partial_k p=k^\sharp\ne0$ on the
nonzero null cone. Off that cone the leading symbol is invertible, giving
the elliptic regularity assertion. On the cone the Hamilton equations are
\[
 \dot x^a=g^{ab}k_b,\qquad
 \dot k_a=-\tfrac12(\partial_a g^{bc})k_b k_c
          =\tfrac12(\partial_a g_{bc})\dot x^b\dot x^c .
\]
The last expression equals $\Gamma^c{}_{ab}k_c\dot x^b$ for the
Levi--Civita connection. Thus $k$ is parallel along $x$ and $x$ is a
null geodesic. The theorem transports a singularity \emph{if it is
present}; it neither puts one on every ray nor chooses its frequency
sign. At a source singularity the hypothesis on $Pu$ must be rechecked.

In particular, if $Pu$ is smooth, every non-characteristic hypersurface
passes the pullback test for $u$; so do its spacetime derivatives, which
cannot enlarge the wavefront set. Both spacelike and timelike
hypersurfaces are non-characteristic for this operator. This does not
make a timelike surface a Cauchy surface. Conversely a perfectly smooth
solution has a restriction even to a characteristic surface. These are
different questions from well-posedness for data on that surface.

\begin{exercise}
In two-dimensional Minkowski space, take $u(t,x)=b_-(t-x)$. Check that
$\Box u=0$ distributionally, determine its wavefront set, and follow one
singular covector along a Hamilton curve.
\end{exercise}
\noindent\emph{Hint.} Use the coordinates $(t-x,x)$ to identify the
covectors as $a\,\mathrm{d}(t-x)$ with $a<0$ over $t=x$. The two second
derivatives in $-\partial_t^2+\partial_x^2$ cancel. The metric dual of
that covector points along the line, not across it.

\section{The frequency selected by a Hadamard state}
Return to a smooth time-oriented globally hyperbolic spacetime without
boundary. For real test functions $f,h$, let $W(f,h)=\Lambda([f],[h])$
be the ordered moment of Proposition~\ref{f:gaussian-positivity}, and
extend $W$ bilinearly to complex tests. Assume continuity on spacetime
tests so the Schwartz kernel theorem gives a distribution $W$ on the
product spacetime \cite{HormI}. This $W$ is not a Weyl generator.
With $\hbar=1$ the free equation and CCR read
$P_x W=P_{x'} W=0$ and $W-W^T=iE$, with $E=G_R-G_A$ and
$W^T(x,x')=W(x',x)$. Hermiticity and positivity mean that $W(\bar f,h)$
is a Hermitian form with $W(\bar f,f)\ge0$, using the spacetime volume
density. Which of the two frequency
orientations is its short-distance singularity allowed to carry?

With signature $(-,+,\ldots,+)$ call a causal covector $k$ future if
$k^\sharp$ is future, equivalently $k(v)<0$ on future timelike vectors
$v$ for nonzero causal $k$. A flat positive-energy phase
$e^{-i\omega t+i\mathbf p\cdot\mathbf x}$ has $k_t=-\omega$, not
$+\omega$, for our Fourier transform. Individual plane waves are smooth;
it is the superposition up to unbounded momentum that produces the
singularity. In a massive field the large-momentum directions still
approach the null cone, in agreement with the principal-symbol test.

\begin{theorem}[Microlocal characterization of Hadamard states \cite{Rad,SV}]
\label{f:wf-hadamard}
For a quasifree Klein--Gordon state with distributional two-point function
on a smooth time-oriented globally hyperbolic spacetime without boundary,
the Hadamard condition in four dimensions \cite{Rad}, and more generally
in spacetime dimension at least three \cite{SV}, is equivalently
\begin{equation}
 \operatorname{WF}(W)=C^+:=
 \{(x,k;x',-k'):(x,k)\sim(x',k'),\ k^\sharp\text{ future null}\}.
 \label{f:wf-spectrum}
\end{equation}
The relation means cotangents to the same null geodesic, parallel
transported along it, with the null tangent proportional to their metric
duals. It includes $x=x'$ with $k=k'\ne0$ null. It does not require
$x$ to be later than $x'$.
\end{theorem}

In four spacetime dimensions one can see exactly what is being prescribed
in position space. Work in a convex normal neighborhood, taking the pair
in a smaller causally convex neighborhood so no returning null ray leaves
the chart. Let $\gamma$ be the unique local geodesic from $x'$ to $x$,
affinely parametrized from zero to one. Write $\sigma=\sigma_g(x,x')$ for
Synge's world function \eqref{f:world-function}: it is half the signed
squared interval, $\tfrac12\int_0^1 g(\dot\gamma,\dot\gamma)\,\mathrm{d}s$.
For a smooth temporal function $T$ increasing to the future, in length units,
put $\sigma_\epsilon=\sigma+i\epsilon(T(x)-T(x'))+\epsilon^2/2$.
The local singular part for a canonically normalized scalar is
\begin{equation}
 H_\ell=\frac1{8\pi^2}\lim_{\epsilon\downarrow0}
 \left(\frac{U}{\sigma_\epsilon}
       +V\log\frac{\sigma_\epsilon}{\ell^2}\right),\qquad \ell>0 ,
 \label{f:wf-parametrix}
\end{equation}
with the logarithm real on the positive axis and $\ell$ a length. The
coefficients are fixed geometrically up to the smooth freedom of a
parametrix. For instance, off the cone use
$\nabla_a\sigma\nabla^a\sigma=2\sigma$ to calculate
\[
 P(U/\sigma)=\frac{PU}{\sigma}
  +\frac{(4-\Box_g\sigma)U-2\nabla^a\sigma\nabla_a U}{\sigma^2} .
\]
Canceling the second numerator transports $U$ along the geodesic, from
$U(x',x')=1$; the logarithmic term cancels the subsequent singular orders.
On smooth metrics this is a parametrix modulo smooth terms, not a claim
that a series for $V$ converges. At any specified finite differentiability
order one can use sufficiently many transport coefficients
\cite{WaldQFT,HWick}. In the massless flat vacuum $U=1,V=0$. Write
$\tau=t-t'$, $\mathbf z=\mathbf x-\mathbf x'$ and $r=|\mathbf z|$.
For $p=|\mathbf p|$, the angular average of $e^{i\mathbf p\cdot\mathbf z}$
is $\sin(pr)/(pr)$. Thus the damped positive-frequency integral is
\begin{align*}
 \int \frac{\mathrm{d}^3\mathbf p}{(2\pi)^3\,2p}
       e^{-(\epsilon+i\tau)p+i\mathbf p\cdot\mathbf z}
 &=\frac1{4\pi^2 r}\int_0^\infty
         e^{-(\epsilon+i\tau)p}\sin(rp)\,\mathrm{d}p\\
 &=\frac1{4\pi^2[r^2-(\tau-i\epsilon)^2]} .
\end{align*}
Taking $\epsilon\downarrow0$ checks the factor of two from the world
function. It also recovers the Fourier test: near a noncoincident null
pair use the normal coordinate $q=r^2-\tau^2$. For $\tau>0$ the boundary
value is $b_+(q)$ up to the prefactor, and for $\tau<0$ it is $b_-(q)$.
In either case a singular covector $a\,\mathrm{d}_x q$ in the first
argument has $k_t=-2a\tau<0$ by \eqref{f:wf-bv}.
The local Hadamard requirement is that $W-H_\ell$ be smooth, using the
full parametrix modulo smooth terms. A change of $\ell$ adds a smooth
multiple of $V$, not a new frequency direction. $H_\ell$ alone
need not be a positive bisolution. A true boundary may introduce waves
reflected along additional rays; the bulk theorem does not supply a
global wavefront condition for an unspecified boundary problem.

\Needspace{.47\textheight}
Two consequences separate measurement from subtraction. First, if $W_1$
and $W_2$ obey \eqref{f:wf-spectrum} and have the same commutator, their
difference is symmetric and has wavefront set contained in $C^+$. Its
transpose instead has past first covector. Since these cones are disjoint,
the difference is a smooth symmetric bisolution. Smooth does not mean
zero: adding $c\,v(x)v(x')$ with $c\ge0$ and $v$ a smooth real solution
preserves positivity and the commutator and can change a fluctuation.
Second, restricting both arguments of $W$ to a spacelike hypersurface is
permitted: its nonzero conormals are timelike and cannot be either null
entry in \eqref{f:wf-spectrum}. On a null hypersurface two arguments on
the same generator can instead have a wavefront pair normal in both
slots. Restricting to $x=x'$ is yet another
map: $(k,-k)$ is a diagonal conormal, so the ordinary pullback test
fails. By contrast $W^2$ as a distribution of \emph{two} arguments passes
the product test because the first covectors cannot be opposite.

\begin{exercise}
Restrict both arguments of the four-dimensional flat massless vacuum to
one inertial clock $\gamma(t)=(t,\mathbf0)$. Why is the pullback allowed
even at equal clock readings? As a distribution of $\tau=t-t'$, compute
the Fourier transform of the result and check its spectral sign.
\end{exercise}
\noindent\emph{Solution.} A nonzero null covector cannot annihilate the
clock's timelike tangent. The restricted distribution is
$-(4\pi^2)^{-1}(\tau-i0)^{-2}=(4\pi^2)^{-1}\partial_\tau b_-$, as one
also finds by doing the exponentially damped mode integral at $r=0$.
Equation~\eqref{f:wf-bv} gives
$-\xi\,\mathbf 1_{\xi<0}/(2\pi)$, a nonnegative spectral density on the
negative temporal Fourier ray. This is a distribution on two clock
readings, not a finite pointwise variance at one reading.
